\chapter*{TFPT und Universalraum\\Update-Paper v1.4}
\noindent\textsf{14. September 2026 · Änderungen gegenüber dem vollständigen Hauptbuch v1.3}

Dieses Update fasst die Prüfung der fünf neuesten Berichte und die danach eigenständig gerechnete Forschungsrunde zusammen. Das zugehörige Hauptbuch v1.4 erhält alle bisherigen Grundlagen, Herleitungen und Abbildungen. Es enthält zusätzlich die vollständigen neuen Argumente, die T1--T8-Zuordnung und die Ressourcen- und Quellengrenzen.

\begin{bild}{Die wichtigste Änderung, einfach erklärt}
Aus einer Sammlung passender Bauteile ist ein vollständig durchgerechnetes kleines Experiment geworden: Zustand herstellen, aufzeichnen, zurückentwickeln, Ergebnis prüfen. Wir wissen jetzt auch, wie viele erfolgreiche Versuche dabei übrig bleiben. Noch fehlt die Herleitung, warum der ursprüngliche Compiler genau diese Bauteile und Bedienmöglichkeiten erzeugt. Der neue Fortschritt liegt im kontrollierten Experiment, nicht in einer bereits fertigen Theorie des Universums.
\end{bild}

\section*{1. Was an den fünf neuen Berichten trägt -- und was nicht}
Der gewichtete Bandbeweis, der Feedbackattraktor und die mikroskopische Sternlösung tragen im deklarierten Modell. Die Standard-$E_8$-Phasenrechnung und die lokalen F4-Zahlen lassen sich unabhängig reproduzieren. Drei Aussagen mussten enger gefasst werden:
\begin{itemize}
\item Dichte Gleitkomma-Diagonalisierung liefert keinen exakten Vierfachheitsbeweis. Die Vierfachstruktur ist numerisch gut gestützt, aber nicht abschließend zertifiziert.
\item Eine Tabelle aller 64 Sektoren enthält überwiegend nur niedrige \emph{nackte} Eigenmoden und darauf ausgewertete Korrekturen. Das ist keine vollständige Neudiagonalisierung des korrigierten Operators.
\item Lokalität unterscheidet nicht nur eine globale Bank von Vermittlern pro Kante. Eine Bank \emph{pro Zelle} ist ebenfalls lokal. Deshalb wählt Lokalität allein weder die Architektur noch das Vorzeichen der Gapkorrektur.
\end{itemize}
Die fünf Originaltexte sind unverändert und mit Prüfsummen im Quellenpaket enthalten. Ihre Ergebnisse werden nicht pauschal als bestätigte Lösungen übernommen.

\section*{2. Eine neue Vereinfachung der Dynamik}
Für kantenlokale Vermittler lässt sich die vierte Ordnung schreiben als
\begin{equation}
F_{4,\mathrm{edge}}=\sum_v A_v(A_v-I)\succeq0,
\qquad A_v=\sum_{e\ni v}(I-S_e).
\end{equation}
Die Ganzzahligkeit des lokalen Spektrums wurde durch eine exakte 720-dimensionale reguläre $S_6$-Rechnung kontrolliert. Zusätzlich gilt auf dem ganzen Materieraum
\begin{equation}
F_{4,\mathrm{edge}}\succeq H_0^2-76H_0+1440I.
\end{equation}
Das verkleinert eine komplizierte Korrektur zu lokalen positiven Bausteinen. Positivität des Operators bedeutet dabei nicht automatisch eine positive \emph{Gapänderung}; Grund- und angeregter Zustand verschieben sich unterschiedlich.

Im neuen 24.024-dimensionalen Singulettlauf wurde anschließend der vollständige trunkierte Operator $H_0+0{,}00125F_{4,\mathrm{edge}}$ im niedrigen Spektralbereich neu diagonalisiert. Bei $t/\Delta=1/20$ finden wir
\begin{equation}
E_0/J\simeq11{,}96050741,\quad E_1/J\simeq12{,}44698494,
\quad\mathrm{Gap}/J\simeq0{,}48647753.
\end{equation}
Das erste gefundene Niveau ist vierfach; die acht berechneten Eigenpaare haben maximale Residualnorm etwa $4{,}75\cdot10^{-14}$. Dies ist weder das gesamte Spektrum noch eine rigorose Fehlerzertifizierung oder die volle Mikrodynamik. Die neue Operatoruntergrenze hilft beim Nichtsingulettvergleich; die dafür verwendeten bisherigen nackten Sektorminima benötigen noch zertifizierte Intervalle.

Eine zweite Verbesserung betrifft den \emph{äußeren} Abstand zum Vermittlerband: Für Vermittler pro Kante beweist ein rationaler Schur-/LDL-Test mindestens $0{,}7\Delta$ bei $|t|/\Delta\le1/20$, gegenüber $0{,}4\Delta$ bei der zellintern geteilten Bank. Ein eigener erster Versuch mit $0{,}8\Delta$ scheiterte und ist als Negativkontrolle dokumentiert. Diese äußere Bandtrennung beweist nicht den inneren Singulettgap und nicht den Vielzellenlimes.

\section*{3. Ein gemeinsamer mikroskopischer Ablauf mit Rohwahrscheinlichkeiten}
Der 544-dimensionale Stern besitzt 14 verschiedene Energien. Mit einem kontrollierten Zeitentwicklungsschritt pro unerwünschtem Niveau liefert
\begin{equation}
P_{\Omega_d}=\prod_{j=1}^{13}
\frac{I+e^{-i\tau_j(H-E_0)/\hbar}}2,
\qquad \tau_j=\frac{\pi\hbar}{E_j-E_0},
\end{equation}
einen exakten postselektierten Grundzustandsfilter im idealen Modell. Er wurde auf der vollständigen Matrix geprüft, einschließlich oberem und dunklem Band. Ein nur auf den unteren Energien getesteter Filter reicht nicht.

Die Projektion auf leere Vermittler verbindet diesen angekleideten Zustand mit dem ursprünglichen Materiezustand:
\begin{equation}
P_{\mathrm{bare}}P_{\Omega_d}P_{\mathrm{bare}}=wP_\Omega,
\qquad w\simeq0{,}9856429312.
\end{equation}
Damit lässt sich eine zusätzliche ideale Entkleidungsoperation vermeiden, allerdings gegen eine ausdrücklich bezahlte Erfolgswahrscheinlichkeit. Aus zwei elementaren Paarsinguletts mit Anfangsüberlapp $1/6$ folgt:
\begin{center}
\begin{tabular}{lr}
\toprule Größe&Unbedingte Wahrscheinlichkeit\\\midrule
Erfolgreiche Präparation&$w^2/6\simeq0{,}16191533$\\
Gesamter Erfolg mit behaltenem Record&$w^4/6\simeq0{,}15729945$\\
Gesamter Erfolg mit frischem Record&$17w^4/192\simeq0{,}08356533$\\\bottomrule
\end{tabular}
\end{center}
Das Verhältnis bleibt exakt $17/32$. Diese Rohwerte umfassen Start- und Endfilter; sie ersetzen nicht stillschweigend ältere Zahlen mit anderem Anfangszustand. Präparation allein braucht im Mittel etwa 6,18 unabhängige Versuche.

Auch die Aufzeichnung wird einfacher: Der resonante Puls liefert direkt $U_{\mathrm{res}}^\dagger Q U_{\mathrm{res}}=R\oplus I$. Für diese Recordoperation sind zusätzliche Viertelphasen nicht nötig; die Inversion des Pulses und die Belegungsabfrage bleiben benötigte Kontrollen.

\begin{grenze}{Kosten und Herkunft bleiben sichtbar}
Ein Start- und Endfilter benötigen zusammen 26 kontrollierte schwache Zeitentwicklungen mit Zeitbudget etwa $6345{,}66\hbar/\Delta$, zusätzlich zu resonanten Records, Belegungsmessungen und Reset. Exakte Spektralzeiten und kontrolliertes $H$ sind deklarierte ideale Ressourcen. Der Ablauf benutzt eine kontrollierte Modellfamilie, nicht einen unverändert frei laufenden Hamiltonoperator. Eine Ableitung all dieser Handgriffe aus dem primitiven TFPT-Compiler fehlt.
\end{grenze}

\section*{4. Weitere eigene Ergebnisse: Phasen und robuste Präparation}
Im kantenlokalen Modell mit erhaltener lokaler Ladung $n_f(v)+\sum_{e\ni v}n_e=1$ bilden besetzte Vermittler eine Paarung. Die fermionischen Vorzeichen lassen sich auf diesem erreichbaren Sektor in allen Ordnungen durch die Diagonale
\begin{equation}
D(M)=(-1)^{\text{Anzahl der Kreuzungen von }M}
\end{equation}
auf die Tensorbeschreibung abbilden. Der vollständige $K_4$-Test umfasst 940 Zustände und 1584 Erzeugungsübergänge. Hopping, eine gemeinsam benutzte Modenbank und die native Clock-Phase sind nicht automatisch eingeschlossen. Die Standard-$E_8$-Cocycle-Rechnung bleibt außerdem von der Identifikation der ursprünglichen TFPT-Phasen zu unterscheiden.

Für den geprüften Mess-/Reset-Feedbackkanal ist der Zielzustand der einzige stationäre Zustand. Mit $r\simeq0{,}97542071$ genügen 556 Zyklen für Infidelität $10^{-6}$. Neu ist die Fehlerkontrolle: Ein Fehler $\epsilon$ pro Zyklus in halber Diamantnorm ergibt einen asymptotischen Infidelitätsboden höchstens
\begin{equation}
\frac{\epsilon}{1-r}\simeq40{,}68466\epsilon.
\end{equation}
Für 4096 unabhängig gereinigte Zellen genügen ideal 890 Zyklen je Zelle für globale Infidelität $10^{-6}$, auch bei anfangs verschränktem Zustand. Das ist keine Grundzustandspräparation eines wechselwirkenden Vielzellensystems und keine kostenlose Entropieentsorgung.

\section*{5. Was auf T1--T8 konkret weitergearbeitet wurde}
\begin{longtable}{L{10mm}L{140mm}}
\toprule Tor&Neue Rechnung und verbleibende Grenze\\\midrule\endhead
T1&Exakte Rang-Ausnahmen $a=1/7,1/3$; Kontext- und Strahlenentropie wählen verschieden. Ohne Auslesevertrag kein eindeutiges Einfachheitsprinzip.\\
T2&Vakuum-Ladungskoeffizienten bis Niveau 32 und Halbladungs-Paritätsschranke. Die native Halbladung samt Energie- und Adjungiertenkontrolle bleibt offen.\\
T3&Skalierende Wärmetrace-Proben in Dimension eins bis vier. Jede hineingesteckte Dimension erscheint; drei Raumdimensionen und ein gemeinsamer Lichtkegel werden nicht ausgewählt.\\
T4&Ein tatsächlich aufgebauter zweidimensionaler Fluss-Overlap-Operator liefert Index drei bei gewähltem Fluss drei. Kontrollflüsse liefern andere Zahlen; Familie, Maß und Spiegelentkopplung sind nicht hergeleitet.\\
T5&Bessere äußere Bandtrennung, positive lokale vierte Ordnung und vollständiger mikroskopischer Filter. Uniforme wechselwirkende Skalierung und kontrollierter Kontinuumslimes fehlen.\\
T6&Drei Nullmoden geben bei konstantem Skalarüberlapp nur gleiche Massen. Der gemeinsame ACT-Test der einfachen Inflationsbranche zeigt eine nicht durch freie Wahl von $N$ entfernbare Spannung.\\
T7&Acht Weyl-Knoten im einfachen Sinusmodell; ein gapped Tensorbilinear bekommt durch TT-Projektion keinen masselosen Pol. Ein angenommener weicher Spin 2 muss universell koppeln.\\
T8&Ein durchgerechneter gemeinsamer operativer Ablauf liegt vor. Sein vollständiges physikalisches Quellfunktional und die native Herkunft der Kontrollen fehlen.\\\bottomrule
\end{longtable}

Der ACT-Test eliminiert $N$ aus der einfachen Branche und erhält $A_s(1-n_s)^2=c_3^7/(6\pi^2)$. Mit festem $c_3=1/(8\pi)$ und Kalibrierung an $A_s$ folgt $n_s\simeq0{,}96468$, gegenüber $0{,}9752\pm0{,}0030$ in der betrachteten P-ACT-LB2-Spalte der \href{https://arxiv.org/abs/2503.14452v2}{ACT-DR6-Tabelle 5}. Das entspricht etwa 3,51 marginalen Standardabweichungen, nicht einer gemeinsamen Likelihoodauswertung. Die Kalibrierung an $n_s$ ergibt stattdessen die 2,0284-fache Amplitude. Ein gemeinsamer Transfer statt getrennter Zahlenanpassung ist hier der notwendige nächste Schritt.

Für RH und Faktorisierung fehlen aktuell gepinnte Originalquellen; die Prüfung dieser neueren externen Routen bleibt daher blockiert. Es wurden keine Pins oder Beweisstatus geändert. Die vorliegenden physikalischen Rechnungen lösen weder RH, P versus NP noch allgemeine Faktorisierung. Auch dunkle Materie, dunkle Energie, Baryogenese, starkes CP und Schwarze Löcher bleiben an die fehlende vollständige Feld- und Gravitationsdynamik gebunden.

\section*{6. Die nächsten sechs Fragen in einfacher Reihenfolge}
\begin{enumerate}
\item \textbf{Wer bedient das Labor?} Jede verwendete Kontrolle auf konkrete primitive Compileroperationen zurückführen.
\item \textbf{Welche Bauweise ist wirklich vorgegeben?} Zellbank und Kantenmoden anhand der Quelle unterscheiden, nicht anhand des erwünschten Resultats.
\item \textbf{Ist die Vierfachstruktur beweisbar?} Symmetrieprojektoren und zertifizierte Sektoruntergrenzen mit der neuen positiven F4-Form verbinden.
\item \textbf{Funktioniert es mit Fehlern und vielen Zellen?} Zeitbedarf, Rücksetzen und Entropieabfluss in einer gemeinsamen Ressourcenrechnung beherrschen.
\item \textbf{Entsteht eine gemeinsame Welt?} Dieselbe skalierende Quelle muss Raumzeit, chirale Materie und einen dynamischen Spin-2-Sektor tragen.
\item \textbf{Treffen dieselben Parameter mehrere Messungen?} Flavour und Kosmologie gemeinsam mit ihrem Transfer prüfen, ohne für jede Zahl nachzujustieren.
\end{enumerate}

\section*{Prüfumfang dieser Revision}
Drei neue unabhängige leichte Prüfer: 1287, 6939 und 6 Bedingungen, insgesamt 8232; normale und optimierte Ausführung bytegleich. Sieben gezielte Fehlvarianten werden erkannt. Hinzu kommt der separate neue 24.024D-Spektrallauf mit acht niedrigen korrigierten Eigenpaaren. Das sind überwiegend einzelne Zustands- und Übergangsprüfungen, nicht 8232 unabhängige Entdeckungen. Die Großrechnung ist numerisch, die analytischen Argumente sind nicht Lean-formalisiert. Alle vollständigen T1--T8-Tore bleiben offen.

\noindent\textsf{Vollständige Herleitungen, Grundlagen und Quellen: Hauptdokument v1.4. Rohwerte, Eingaben und reproduzierbare Prüfer: zugehöriges Quellenpaket.}
